
Weight space learning methods operate on neural network parameters as data, enabling applications in model analysis and synthesis. Existing approaches succeed within single architecture families but have not been extended to heterogeneous populations. We test whether Deep Sets encoders augmented with architecture embeddings and MSE-based equivariance loss can learn shared representations across CNNs, Transformers, and RNNs. We train on a synthetic model zoo of 1000 models with 1000-dimensional weight vectors and evaluate using three metrics: reconstruction error (information preservation), frozen-K generalization (architecture independence), and kernel robustness (permutation invariance). All three metrics fail: reconstruction error reaches 19.18\% (target <10\%), frozen-K generalization reaches 10.31\% (target <10\%), and kernel robustness achieves 0.00\% (target ≥90\%). The kernel robustness result indicates complete failure to learn permutation invariance despite explicit equivariance loss training. We identify three contributing factors: MSE-based equivariance loss does not enforce group structure in this configuration, architecture embeddings may anchor representations to family-specific coordinates, and K=32 quotient dimensions are insufficient for the tested weight dimensionality. These findings document a specific failed configuration and suggest alternative research directions including contrastive learning for equivariance, architecture-agnostic encoding, and validation on homogeneous populations before attempting cross-architecture extension.
